% ---------------------------------------------------------------------------
% Resumen (ES) + Abstract (EN) + Palabras clave / Keywords.
% Draft only: the project is at its start, so no results are claimed.
% ---------------------------------------------------------------------------

\chapter*{Resumen}
\addcontentsline{toc}{chapter}{Resumen}
\markboth{Resumen}{Resumen}

Este Trabajo de Fin de Grado aborda el desarrollo de un dashboard inteligente
para la gestión basada en datos aplicada al rendimiento de corredores. En el
ámbito deportivo, el uso de datos ha adquirido una relevancia creciente para
mejorar el rendimiento, prevenir lesiones y optimizar los entrenamientos. Sin
embargo, buena parte de las plataformas y dispositivos actuales ofrecen
interfaces cerradas o de uso limitado para el análisis personalizado, lo que
dificulta convertir el volumen de datos generado en decisiones concretas.

El proyecto propone una solución que integra métricas derivadas de fuentes de
datos deportivas ---como el ritmo, la frecuencia cardíaca, la distancia, la
potencia o la cadencia--- en un panel interactivo orientado tanto a deportistas
individuales como a entrenadores y equipos. Sobre esa base, se plantea un modelo
de análisis que permita detectar patrones de evolución y un módulo predictivo
basado en regresión lineal y modelos de series temporales para estimar la
progresión del rendimiento a corto y medio plazo. El trabajo se enmarca en las
metodologías de \emph{Data-Driven Management} aplicadas al contexto deportivo.
Dado que el proyecto se encuentra en su fase inicial, este documento describe el
problema, el enfoque propuesto y la contribución esperada, sin anticipar
resultados.

\vspace{0.5\baselineskip}
\noindent\textbf{Palabras clave:} análisis de datos deportivos, rendimiento de
corredores, dashboard interactivo, aprendizaje automático, series temporales,
Strava.

\chapter*{Abstract}
\addcontentsline{toc}{chapter}{Abstract}
\markboth{Abstract}{Abstract}

This Bachelor's Thesis addresses the development of a smart dashboard for
data-driven management applied to runner performance. In the sports domain, the
use of data has become increasingly relevant to improve performance, prevent
injuries and optimise training. However, many current platforms and devices
offer closed or limited interfaces for customised analysis, which makes it
difficult to turn the volume of generated data into concrete decisions.

The project proposes a solution that integrates metrics derived from sports data
sources ---such as pace, heart rate, distance, power and cadence--- into an
interactive panel aimed at both individual athletes and coaches or teams. On
that basis, it outlines an analysis model able to detect patterns of evolution
and a predictive module based on linear regression and time-series models to
estimate performance progression in the short and medium term. The work is
framed within Data-Driven Management methodologies applied to the sports
context. Since the project is at an early stage, this document describes the
problem, the proposed approach and the expected contribution, without
anticipating results.

\vspace{0.5\baselineskip}
\noindent\textbf{Keywords:} sports data analytics, runner performance,
interactive dashboard, machine learning, time series, Strava.
